\section*{Resumen Ejecutivo y Alcance del Informe}
\addcontentsline{toc}{section}{Resumen Ejecutivo y Alcance del Informe}

El presente informe desarrolla y amplía el análisis urbanístico correspondiente a la \textbf{Tarea 4} sobre un sector residencial urbanizable de $6\text{ ha}$ ($S_{\text{bruta}} = 60.000\text{ m}^2\text{s}$, determinado conforme a la última cifra del DNI del alumno, $4 + 2 = 6\text{ ha}$). Partiendo del balance cuantitativo de los cuatro escenarios base ($S_1$--$S_4$) y de la restricción morfológica de altura máxima en \textbf{dos plantas} ($N_p = 2$), el estudio integra el marco jurídico canario y estatal (\emph{Ley 4/2017}, \emph{Decreto 183/2018} y \emph{RDL 7/2015}), la economía empírica del promotor inmobiliario basada en cuartiles de mercado y costes físicos reales, la reserva de vivienda protegida (VPO) heredada de la \emph{Tarea 3}, y un contraste territorial empírico frente al Plan General de Ordenación de Santa Cruz de Tenerife (\emph{El Tablero} y \emph{Llano del Moro 3}).

\subsection*{Conclusiones Principales}

\begin{enumerate}
  \item \textbf{Imposibilidad Geométrica de $C_e = 1,0$ en Dos Plantas (Sección~\ref{sec:modelo_normativo}):} Al descontar un viario fijo del $10\%$ ($6.000\text{ m}^2\text{s}$) y cesiones dotacionales acopladas a la edificabilidad ($0,45\text{ m}^2\text{s}/\text{m}^2\text{c}$), agotar $60.000\text{ m}^2\text{c}$ en dos plantas exigiría una ocupación de parcela (\emph{Building Coverage Ratio}, $\text{BCR}$) del $111,11\%$, físicamente imposible. Bajo un límite urbanístico realista del $\text{BCR} = 90\%$, el máximo aprovechamiento materializable se sitúa en $\text{Gross FAR} = C_e^* = 0,8950\text{ m}^2\text{c}/\text{m}^2\text{s}$ ($53.702\text{ m}^2\text{c}$, $42$ parcelas).
  \item \textbf{Salidas Jurídicas Reales y Cesión Municipal (Sección~\ref{sec:realidad_juridica}):} Los garajes obligatorios bajo rasante declarados no computables no absorben el déficit de edificabilidad de $S_1$ ($-19\%$) o $S_2$ ($-55\%$). Las verdaderas vías jurídicas para materializar dicho techo son la \emph{Transferencia de Aprovechamiento Urbanístico} (\emph{TAU}), los espacios \emph{bajo cubierta} computables o los semisótanos/plazas excedentarias donde la ordenanza los compute. Asimismo, la cesión municipal del $10\%$ (\emph{PMS}) se entrega en parcelas netas urbanizadas libres de cargas, asumiendo el promotor el $100\%$ de la urbanización del sector.
  \item \textbf{Óptimo Económico frente al Límite Normativo (Sección~\ref{sec:mercado_economico}):} La valoración por cuartiles de mercado ($Q_{25}$--$Q_{75}$) revela asimetría positiva en baja densidad y compresión de márgenes al densificar. Sumado al escalón de coste físico del garaje subterráneo al superar el $70\%$ de ocupación ($\text{BCR} > 70\%$), el margen neto esperado $\mathbb{E}[\Pi]$ alcanza su máximo en el interior de la frontera ($\text{Gross FAR}^* \approx 0,58$--$0,61$, $\text{BCR} \approx 45\%$--$49\%$), muy por debajo del máximo edificable ($0,895$).
  \item \textbf{Desacoplamiento por Reserva del $40\%$ de VPO (Sección~\ref{sec:vpo_mixta}):} Al imponer la reserva del $40\%$ sobre el techo edificable total ($\text{m}^2\text{c}$), incluyendo el $10\%$ municipal en VPO, la ocupación se desdobla en $(\text{BCR}_M, \text{BCR}_V)$. Incluso imputando un sobrecoste unificado de garaje comunitario en sótano ($+250\text{ EUR}/\text{m}^2\text{c}$) al compactar la VPO, concentrar la vivienda protegida al $90\%$ de $\text{BCR}$ en $2$ plantas ($\text{Gross FAR}^* = 0,670$, $44,7\text{ M EUR}$) o en torres de $5$ plantas ($\text{Gross FAR}^* = 0,715$, $47,1\text{ M EUR}$) protege el suelo de baja densidad de mercado libre.
  \item \textbf{Contraste Teoría frente a Realidad en Santa Cruz de Tenerife (Sección~\ref{sec:teoria_realidad}):} La auditoría sobre \emph{UA-1 (SSU 5.13 El Tablero, $6,83\text{ ha}$)} y \emph{SSU 5.20 (Llano del Moro 3, $6,03\text{ ha}$)} demuestra que el viario real ($22\%$--$29\%$) y las cesiones de $\text{EL/D/I}$ ($38\%$--$40\%$) superan ampliamente al $10\%$ y $32,2\%$ del modelo teórico, no por mera orografía (pendientes moderadas del $1\%$--$12\%$), sino por la absorción de \emph{Sistemas Generales}, bulevares peatonales, bordes de barranco y déficits dotacionales del barrio colindante. Paralelamente, la parálisis judicial del PGO de 2013 (\emph{Sentencia TSJC 116/2014}) y de la Ordenanza de 2022 (\emph{TSJC 000070/2023}) mantiene solares ya urbanizados de \emph{UA-1} aún sin edificar tras expirar los plazos en 2023.
\end{enumerate}

\begin{tcolorbox}[jumpbox, title={\textbf{Nota Terminológica Internacional}}]
En las gráficas técnicas adjuntas (\texttt{img/}) se emplea la nomenclatura urbanística internacional equivalente a la normativa española: \textbf{BCR} (\emph{Building Coverage Ratio} / \emph{Lot Coverage}) corresponde a la \textbf{ocupación de parcela} ($C_{\text{ocup}}$); \textbf{Gross FAR} al coeficiente de \textbf{edificabilidad bruta} ($C_e = E_{\text{teórica}}/S_{\text{bruta}}$); y \textbf{Net FAR} a la \textbf{edificabilidad neta} sobre parcela lucrativa ($E_{\text{teórica}}/S_{\text{neta}} = \text{BCR} \cdot N_p$). El entorno de verificación cartográfica y cálculo reproducibles ha contado con asistencia de \emph{Google Antigravity} \cite{google_antigravity_2026}.
\end{tcolorbox}
